\documentclass[12pt]{article}

\usepackage{sbc-template}
\usepackage{graphicx,url}
\usepackage[utf8]{inputenc}
\usepackage[brazil]{babel}
\usepackage{booktabs}
\usepackage{tabularx}
\usepackage{tikz}
\usetikzlibrary{shapes.geometric, arrows.meta, positioning}

\sloppy

\title{Análise da Distribuição de Estabelecimentos de Saúde na Região Sul
  e Impactos na Mortalidade:\\Uma Arquitetura de Dados Híbrida Batch/Stream}

\author{Jade S. Hatanaka Marques\inst{1}, Luís Dias Ferreira Soares\inst{1}} 

\address{Instituto de Informática -- Universidade Federal do Rio Grande do Sul (UFRGS)\\
  Porto Alegre -- RS -- Brasil
  \email{jade.hatanaka@inf.ufrgs.br, ldiass@live.com}
}

\begin{document}

\maketitle

\begin{resumo}
A integração de dados de saúde pública no Brasil enfrenta o desafio de
transformar registros administrativos fragmentados e heterogêneos em
inteligência acionável para gestores e profissionais de saúde. Este trabalho
propõe uma plataforma de análise que correlaciona a distribuição de
estabelecimentos de saúde na Região Sul do Brasil — Paraná, Santa Catarina e
Rio Grande do Sul — com desfechos de mortalidade hospitalar, por meio de uma
arquitetura híbrida de processamento em lote (batch) e em tempo real (stream).
O objetivo geral é integrar três fontes públicas — o Sistema de Informações
Hospitalares (SIH/DataSUS), o Cadastro Nacional de Estabelecimentos de Saúde
(CNES) e indicadores socioeconômicos do IBGE — em um banco de dados relacional
estruturado segundo o padrão Medallion (camadas \textit{raw}, \textit{trusted}
e \textit{refined}), habilitando consultas analíticas complexas de
rastreabilidade do paciente e identificação de regiões com alta mortalidade e
baixa oferta assistencial, denominadas ``Silêncio Epidemiológico''. O ambiente
é conteinerizado via Docker, com processamento batch em Apache Spark/PySpark,
simulação de triagem via Apache Kafka e armazenamento em PostgreSQL 16. O
trabalho será avaliado por demonstração ao vivo do pipeline completo e exibição
de mapas coropléticos interativos. A principal contribuição é uma metodologia
replicável de integração de dados públicos de saúde que suporta tanto análise
histórica quanto decisão assistencial em tempo real.
\end{resumo}


%=============================================================================
\section{Introdução}
%=============================================================================

A saúde pública brasileira opera sobre um vasto ecossistema de sistemas de
informação geridos pelo Ministério da Saúde, cujos microdados são disponibilizados
abertamente pelo DataSUS \cite{datasus2024}. Apesar do volume e da riqueza dessas
bases, a utilização efetiva desses dados para suporte à decisão clínica e
epidemiológica permanece limitada pela fragmentação entre sistemas, pela
heterogeneidade de formatos e pela ausência de ferramentas integradas de análise
\cite{lima2019}. O resultado é um ciclo em que gestores tomam decisões de
alocação de recursos sem visibilidade sobre o fluxo real de pacientes e seus
desfechos.

Um exemplo concreto do problema ocorre na regulação assistencial: um paciente
com quadro clínico grave em um município de pequeno porte pode aguardar horas
por uma transferência regulada enquanto leitos de Unidade de Terapia Intensiva
(UTI) permanecem disponíveis em hospitais próximos, simplesmente porque não há
sistema capaz de cruzar, em tempo real, a demanda de triagem com a oferta de
leitos do Sistema Único de Saúde (SUS). Estima-se que essa ``Janela de Risco''
— o intervalo entre a solicitação de regulação e a efetivação da transferência
— contribui para desfechos adversos evitáveis \cite{mendes2011}.

Este trabalho é adequado ao contexto da Engenharia e Ciência de Dados
aplicada à saúde pública, com foco na Região Sul do Brasil, que possui boa
cobertura e confiabilidade nos registros do SIH e do CNES, além de dados
socioeconômicos do IBGE bem estruturados por setor censitário \cite{ibge2023}.
O objetivo geral é construir uma plataforma capaz de correlacionar latência
na regulação assistencial com desfechos de mortalidade, provendo uma camada
de decisão em tempo real para otimização do fluxo de pacientes críticos.

As tecnologias utilizadas incluem: Python 3.11 com PySpark para pipelines de
transformação em lote \cite{spark2016}; Apache Kafka para ingestão de eventos
de triagem simulados \cite{kafka2021}; PostgreSQL 16 como banco de dados
relacional principal; e Docker Compose para conteinerização e portabilidade
da stack completa. A arquitetura de armazenamento segue o padrão Medallion,
amplamente adotado em plataformas modernas de dados \cite{databricks2021}.

A contribuição esperada é dupla: científica, com uma metodologia de integração
de bases heterogêneas de saúde pública reutilizável em outros contextos
regionais; e social, com um protótipo funcional que demonstra como dados já
existentes podem subsidiar decisões de gestão assistencial sem custo adicional
de coleta.

O restante do artigo está organizado da seguinte forma: a Seção~2 descreve
detalhadamente as bases de dados utilizadas; a Seção~3 apresenta o ambiente
de desenvolvimento; a Seção~4 detalha a metodologia adotada; a Seção~5
especifica as funcionalidades do sistema; a Seção~6 apresenta o cronograma de
desenvolvimento; e a Seção~7 traz as considerações finais.


%=============================================================================
\section{Descrição Detalhada da Base de Dados}
%=============================================================================

O ecossistema de dados do projeto é composto por três fontes públicas
integradas pela chave de município no padrão IBGE de seis dígitos.

\subsection{SIH/DataSUS — Sistema de Informações Hospitalares}

O SIH registra todas as internações financiadas pelo SUS por meio das
Autorizações de Internação Hospitalar (AIH). Os arquivos são publicados
mensalmente no FTP do DataSUS no formato proprietário \texttt{.dbc}
(compressão sobre DBF), identificados pelo padrão
\texttt{RD\{UF\}\{AAAAMM\}.dbc}. Para a Região Sul, o volume estimado é de
aproximadamente 500\,mil registros por estado por ano, totalizando cerca de
1,5\,milhão de registros anuais considerando PR, SC e RS no período 2022--2024.

Os campos centrais para a análise são apresentados na Tabela~\ref{tab:sih}.

\begin{table}[ht]
\centering
\caption{Campos principais do SIH/DataSUS utilizados na análise}
\label{tab:sih}
\begin{tabularx}{\textwidth}{llX}
\toprule
\textbf{Campo} & \textbf{Tipo} & \textbf{Descrição} \\
\midrule
\texttt{N\_AIH}     & VARCHAR(13) & Número da AIH — dado sensível, armazenado apenas como hash SHA-256 \\
\texttt{MUNIC\_MOV} & CHAR(6)     & Município de atendimento (código IBGE 6 dígitos) \\
\texttt{MUNIC\_RES} & CHAR(6)     & Município de residência — diferença indica transferência \\
\texttt{DT\_INTER}  & DATE        & Data de internação \\
\texttt{DT\_SAIDA}  & DATE        & Data de saída \\
\texttt{MORTE}      & SMALLINT    & Indicador de óbito: 0 = não, 1 = sim \\
\texttt{DIAG\_PRINC}& VARCHAR(4)  & CID-10 do diagnóstico principal \\
\texttt{CGC\_HOSP}  & VARCHAR(14) & CNPJ do hospital (FK implícita para o CNES) \\
\bottomrule
\end{tabularx}
\end{table}

\noindent\textbf{Regras de qualidade:} registros com \texttt{DT\_SAIDA < DT\_INTER}
são descartados; \texttt{SEXO} com valores 0 ou 9 são convertidos para nulo;
datas de nascimento inválidas (\texttt{00000000}) são tratadas como nulo;
registros com \texttt{MORTE = 1} e \texttt{DT\_SAIDA} nulo são descartados.

\subsection{CNES — Cadastro Nacional de Estabelecimentos de Saúde}

O CNES registra todos os estabelecimentos de saúde do país com informações
de infraestrutura, equipamentos e leitos. São utilizados dois grupos de
arquivos mensais: \texttt{ST\{UF\}AAAAMM.dbc} (dados gerais do estabelecimento)
e \texttt{LT\{UF\}AAAAMM.dbc} (leitos por tipo). Para a Região Sul, a
competência de dezembro de 2024 contém aproximadamente 12\,mil estabelecimentos
ativos. O campo \texttt{QT\_SUS} (quantidade de leitos disponível ao SUS) é
o indicador central para o cálculo de disponibilidade assistencial.

\subsection{IBGE — Indicadores Socioeconômicos}

Os dados do Censo Demográfico 2022 são utilizados em nível de setor censitário,
fornecendo população total, domicílios e estimativa de renda média. Para a
Região Sul, o Censo 2022 conta com aproximadamente 47\,mil setores censitários.
Um ponto crítico de integração é a divergência de código de município: o IBGE
publica códigos de 7 dígitos (com dígito verificador), enquanto o SIH e o CNES
utilizam 6 dígitos. A integração exige a truncagem:
\texttt{SUBSTRING(ibge.CD\_GEOCODM, 1, 6) = sih.MUNIC\_MOV}.

\subsection{Modelagem do Banco de Dados}

O banco de dados PostgreSQL 16 adota o padrão Medallion com três schemas:
\textbf{raw} (espelho das fontes originais, mínima transformação),
\textbf{trusted} (dados limpos, tipados e com anonimização) e
\textbf{refined} (agregações analíticas para suporte à decisão). A
rastreabilidade do paciente é garantida pelo armazenamento do hash
SHA-256 do campo \texttt{N\_AIH}, em conformidade com a LGPD
(Lei nº 13.709/2018, Art. 11), que classifica dados de saúde como sensíveis.
O volume total estimado após carga inicial é de aproximadamente 4,5\,milhões
de registros entre as tabelas \texttt{trusted.internacoes} e
\texttt{raw.sih\_internacoes}.


%=============================================================================
\section{Ambiente de Desenvolvimento}
%=============================================================================

Todo o ambiente é executado localmente via Docker Compose, garantindo
portabilidade e reprodutibilidade independentemente do sistema operacional
do desenvolvedor. A Tabela~\ref{tab:stack} resume os componentes da stack.

\begin{table}[ht]
\centering
\caption{Stack tecnológica do projeto}
\label{tab:stack}
\begin{tabularx}{\textwidth}{llX}
\toprule
\textbf{Componente} & \textbf{Versão} & \textbf{Função} \\
\midrule
Docker Compose      & 3.x             & Orquestração de todos os serviços em containers isolados \\
PostgreSQL          & 16              & SGBD relacional principal — schemas raw/trusted/refined \\
Apache Spark        & 3.5 (Bitnami)   & Processamento batch distribuído — Master + Worker \\
Apache Kafka        & 7.6.1 (Confluent)& Broker de mensagens para pipeline de eventos de triagem \\
Apache Zookeeper    & 7.6.1           & Coordenação do cluster Kafka \\
Kafka UI            & latest          & Interface web para inspeção de tópicos e mensagens \\
Python              & 3.11            & Linguagem principal dos pipelines ETL e stream \\
PySpark             & 3.5.x           & API Python para processamento distribuído no Spark \\
pysus               & latest          & Leitura de arquivos .dbc do DataSUS em Python \\
kafka-python        & 2.0.x           & Producer/Consumer Kafka em Python \\
psycopg2            & 2.9.x           & Driver PostgreSQL para Python \\
geopandas / Folium  & 0.14 / latest   & Análise espacial e mapas coropléticos \\
Plotly              & latest          & Gráficos interativos para o dashboard \\
\bottomrule
\end{tabularx}
\end{table}

\noindent\textbf{Justificativa das escolhas:}

\textit{PostgreSQL} foi escolhido em detrimento de bancos NoSQL porque o escopo
da disciplina requer um SGBD relacional completo com suporte a DDL, constraints,
schemas, transações e colunas computadas (\texttt{STORED}). Em ambiente de
produção, as camadas \textit{raw} e \textit{trusted} seriam substituídas por
armazenamento em objeto (S3/GCS) e Delta Lake.

\textit{Apache Spark} em modo standalone via Docker oferece o mesmo modelo de
programação distribuído (RDDs, DataFrames, planos de execução físico) que seria
utilizado em um cluster gerenciado (EMR, Dataproc, Databricks), sendo adequado
para fins de aprendizado e portfólio.

\textit{Apache Kafka} é o padrão de mercado para pipelines de eventos e foi
escolhido sobre alternativas mais simples (RabbitMQ, Redis Streams) pela
relevância curricular e pela capacidade de demonstrar particionamento,
grupos de consumidores e retenção configurável de mensagens.

\textit{Docker Compose} elimina dependências de configuração manual e permite
que toda a stack seja iniciada com um único comando (\texttt{docker compose up -d}),
o que é fundamental para reprodutibilidade e para a demonstração final.


%=============================================================================
\section{Metodologia}
%=============================================================================

O desenvolvimento segue uma abordagem incremental em seis fases sequenciais,
onde cada fase produz um artefato testável antes de avançar para a próxima.
A Figura~\ref{fig:fluxo} apresenta o fluxograma das etapas.

\begin{figure}[ht]
\centering
\begin{tikzpicture}[
  node distance=0.6cm and 0.3cm,
  box/.style={rectangle, rounded corners=4pt, draw, fill=blue!10,
              text width=3.2cm, align=center, minimum height=0.8cm, font=\small},
  arrow/.style={-{Stealth}, thick}
]
  \node[box] (col)  {Coleta de Dados\\(FTP DataSUS/IBGE)};
  \node[box, below=of col] (conv) {Conversão\\.dbc $\to$ Parquet};
  \node[box, below=of conv] (etl)  {ETL Batch\\(PySpark)};
  \node[box, below=of etl]  (pg)   {Carga PostgreSQL\\raw/trusted/refined};
  \node[box, right=1.8cm of pg]   (kafka) {Simulação Stream\\(Kafka Producer)};
  \node[box, above=of kafka] (cons) {Motor de Decisão\\(Kafka Consumer)};
  \node[box, above=of cons]  (esp)  {Análise Espacial\\(Notebooks)};
  \node[box, above=of esp]   (dash) {Dashboard\\Interativo};

  \draw[arrow] (col)   -- (conv);
  \draw[arrow] (conv)  -- (etl);
  \draw[arrow] (etl)   -- (pg);
  \draw[arrow] (pg)    -- (kafka);
  \draw[arrow] (kafka) -- (cons);
  \draw[arrow] (cons)  -- (esp);
  \draw[arrow] (esp)   -- (dash);
  \draw[arrow] (pg)    -- ++(2.5,0) |- (esp);
\end{tikzpicture}
\caption{Fluxograma das etapas de desenvolvimento do projeto}
\label{fig:fluxo}
\end{figure}

\textbf{Fase 1 — Coleta de dados:} os arquivos \texttt{.dbc} são baixados
programaticamente do FTP público do DataSUS usando a biblioteca \textit{pysus}.
Para o SIH, são obtidos os arquivos mensais de PR, SC e RS para o período
2022--2024 (36 arquivos); para o CNES, os arquivos de estabelecimentos (ST) e
leitos (LT) da competência dezembro/2024; para o IBGE, os CSVs de setores
censitários do Censo 2022 para a Região Sul.

\textbf{Fase 2 — Conversão para Parquet:} os arquivos \texttt{.dbc} são
convertidos para o formato Apache Parquet usando \textit{pyarrow}, com
particionamento por UF e competência (ex.: \texttt{data/raw/sih/uf=RS/ano=2024/mes=01/}).
O Parquet é o formato do datalake local, permitindo leitura seletiva por coluna
e compressão eficiente.

\textbf{Fase 3 — ETL Batch (PySpark):} jobs PySpark lêem os Parquets,
aplicam regras de qualidade de dados (filtros de inconsistência, tratamento
de nulos, normalização de tipos), calculam o hash SHA-256 do número de AIH
(anonimização) e calculam colunas derivadas (faixa etária, dias de permanência).
Os dados são carregados no PostgreSQL via JDBC nas camadas \textit{raw} e
\textit{trusted}. A camada \textit{refined} é populada por queries SQL de
agregação que calculam a taxa de mortalidade por município e a disponibilidade
de leitos por 1.000 habitantes.

\textbf{Fase 4 — Simulação de Stream (Kafka):} um \textit{producer} Python
gera eventos JSON representando chegadas de pacientes em pronto-socorro, com
distribuição horária baseada nas estatísticas históricas do SIH. Cada evento
contém: identificador anônimo do paciente, código CNES de origem, CID suspeito
e sinais vitais simulados (FC, PAS, FR, SpO2, Escala de Glasgow). O
\textit{consumer} aplica o Protocolo Manchester de classificação de risco e,
para pacientes classificados como ``Vermelho'' (risco imediato), consulta a
camada \textit{refined} em busca das unidades mais próximas com leitos SUS
disponíveis, gravando o alerta na tabela \texttt{refined.alertas\_triagem}.

\textbf{Fase 5 — Análise Espacial:} notebooks Python utilizam GeoPandas e
Folium para gerar mapas coropléticos de mortalidade por município e mapas de
calor de concentração de estabelecimentos. Uma view materializada
(\texttt{refined.silencio\_epidemiologico}) identifica os municípios acima do
75º percentil de mortalidade e abaixo do 25º percentil de leitos por habitante.

\textbf{Fase 6 — Dashboard e Demonstração:} uma página única interativa em
HTML/Plotly integra os gráficos e mapas gerados. Um \textit{Makefile} com o
alvo \texttt{make demo} sobe toda a stack Docker, carrega uma amostra de dados
e inicia o producer Kafka, permitindo demonstração ao vivo do pipeline em tempo real.

A avaliação do trabalho se dará pela execução da demonstração completa,
verificação da corretude das consultas analíticas sobre os dados carregados e
análise dos mapas gerados.


%=============================================================================
\section{Funcionalidades}
%=============================================================================

As funcionalidades do sistema são divididas em básicas e avançadas, podendo
assumir a forma de consultas SQL, análises estatísticas ou visualizações
interativas.

\subsection{Funcionalidades Básicas}

\textbf{F1 — Distribuição de estabelecimentos por tipo e município:}
consulta sobre \texttt{trusted.estabelecimentos} que retorna, para cada
município da Região Sul, a contagem de unidades por tipo (\textit{TP\_UNID}):
hospitais gerais (04), UPA (07), UBS (02) e centros de especialidade (36).
Permite identificar municípios descobertos de determinado nível de atenção.

\textbf{F2 — Taxa de mortalidade hospitalar por município e competência:}
agrega os registros de \texttt{trusted.internacoes} calculando
$\text{taxa} = (\text{óbitos} / \text{internações}) \times 100$ por município
e mês/ano, populando a tabela \texttt{refined.mortalidade\_por\_municipio},
cujo campo \texttt{taxa\_mortalidade} é uma coluna computada \texttt{STORED}
do PostgreSQL.

\textbf{F3 — Disponibilidade de leitos SUS por 1.000 habitantes:}
a tabela \texttt{refined.disponibilidade\_leitos} consolida o total de leitos
SUS de \texttt{trusted.estabelecimentos} com a população estimada do IBGE,
calculando o indicador $\text{leitos}/1000\text{hab}$ por município. Esse
indicador é utilizado pelo motor de decisão em tempo real e pela análise do
Silêncio Epidemiológico.

\textbf{F4 — Rastreabilidade de internações por estabelecimento:}
consulta parametrizada por CNES que retorna, via hash SHA-256 do n\_AIH, o
histórico anonimizado de internações de um estabelecimento: volume por
competência, taxa de mortalidade interna e proporção de pacientes transferidos
de outros municípios (\texttt{MUNIC\_RES $\neq$ MUNIC\_MOV}).

\subsection{Funcionalidades Avançadas}

\textbf{F5 — Mapa do ``Silêncio Epidemiológico'':}
visualização coroplética interativa (Folium + GeoPandas) que destaca os
municípios identificados pela view materializada
\texttt{refined.silencio\_epidemiologico} — aqueles com taxa de mortalidade
acima do 75º percentil regional e disponibilidade de leitos SUS abaixo do
25º percentil. O mapa é enriquecido com camada de vulnerabilidade
socioeconômica (IDHM e renda média do IBGE), evidenciando a correlação entre
exclusão social e baixa oferta assistencial.

\textbf{F6 — Motor de decisão em tempo real (Kafka):}
ao receber um evento de triagem classificado como ``Vermelho'' (protocolo
Manchester: SpO$_2 < 90\%$ \textbf{ou} Glasgow $< 9$ \textbf{ou} PAS $< 80$\,mmHg),
o consumer consulta as três unidades com maior disponibilidade de leitos SUS
no município de origem e nos municípios vizinhos, publica o alerta no tópico
\texttt{alertas-risco} e registra a recomendação em
\texttt{refined.alertas\_triagem} com \textit{timestamp} para análise
posterior da Janela de Risco.


%=============================================================================
\section{Cronograma}
%=============================================================================

O desenvolvimento está organizado em seis fases ao longo do semestre letivo,
conforme a Tabela~\ref{tab:crono}. Cada fase tem um entregável testável
que valida a conclusão antes do início da próxima.

\begin{table}[ht]
\centering
\caption{Cronograma de desenvolvimento do projeto}
\label{tab:crono}
\begin{tabularx}{\textwidth}{lXll}
\toprule
\textbf{Fase} & \textbf{Atividade} & \textbf{Período} & \textbf{Responsável} \\
\midrule
1a & \texttt{docker-compose.yml}, DDL PostgreSQL (schemas raw/trusted/refined), índices & Abr/1--2 & Jade H. \\
1b & Tópicos Kafka (\texttt{topics.sh}), configuração Zookeeper, CI (\texttt{pyproject.toml}) & Abr/1--2 & Luís D. \\
\midrule
2a & Download SIH: \texttt{download\_sih.py} (PR/SC/RS 2022--2024), conversão .dbc $\to$ Parquet & Abr/3--4 & Jade H. \\
2b & Download CNES: \texttt{download\_cnes.py} (ST + LT), download IBGE: \texttt{download\_ibge.py} & Abr/3--4 & Luís D. \\
\midrule
3a & ETL SIH: \texttt{etl\_sih.py} (limpeza, hash SHA-256, carga raw/trusted); agregações \texttt{aggregate.py} & Mai/1--3 & Jade H. \\
3b & ETL CNES: \texttt{etl\_cnes.py} (join ST+LT, leitos SUS); ETL IBGE: \texttt{etl\_ibge.py}; testes unitários & Mai/1--3 & Luís D. \\
\midrule
4a & Kafka Producer: \texttt{producer.py} (geração de eventos, distribuição horária baseada no SIH) & Mai/4--Jun/1 & Luís D. \\
4b & Kafka Consumer + Protocolo Manchester; DDL \texttt{refined.alertas\_triagem} & Mai/4--Jun/1 & Jade H. \\
\midrule
5a & View materializada \texttt{refined.silencio\_epidemiologico}; análise de transferências inter-municipais & Jun/2--3 & Luís D. \\
5b & Mapas coropléticos (Folium/GeoPandas): mortalidade, leitos, Silêncio Epidemiológico & Jun/2--3 & Jade H. \\
\midrule
6a & Dashboard interativo (Plotly); \texttt{Makefile} com alvo \texttt{make demo}; entrega final & Jun/4--Jul/1 & Jade H. \\
6b & Diagrama ER (\texttt{eralchemy2}); atualização \texttt{pyproject.toml}; documentação final & Jun/4--Jul/1 & Luís D. \\
\bottomrule
\end{tabularx}
\end{table}

As atividades de infraestrutura (Fase 1) já se encontram concluídas, com a
stack Docker operacional e o DDL do banco implementado com os schemas
\textit{raw}, \textit{trusted} e \textit{refined}, seus índices analíticos e
comentários de documentação. As Fases 2 e 3 formam o núcleo crítico do
projeto, pois condicionam todas as análises subsequentes. A divisão de tarefas
busca distribuir uniformemente a carga entre os dois integrantes, com Jade
concentrada no pipeline batch (SIH, ETL, análise espacial) e Luís no pipeline
stream (Kafka, CNES/IBGE, documentação técnica).


%=============================================================================
\section{Considerações Finais}
%=============================================================================

Este trabalho representa uma oportunidade de aplicar, de forma integrada,
os conceitos centrais da disciplina de Projeto de Banco de Dados em um problema
de alto impacto social: a ineficiência na alocação de recursos assistenciais
no SUS. Espera-se consolidar habilidades em modelagem relacional avançada
(schemas, constraints, colunas computadas), processamento distribuído com
Apache Spark, integração de dados heterogêneos e arquitetura de pipelines
híbridos batch/stream.

Do ponto de vista \textbf{científico}, a contribuição é uma metodologia
documentada e replicável para integração das bases públicas SIH, CNES e IBGE
em uma plataforma analítica unificada. A adoção do padrão Medallion e da
anonimização por hash SHA-256 estabelece um modelo que pode ser estendido para
outras regiões do país ou adaptado para outros sistemas de informação em saúde
pública.

Do ponto de vista \textbf{social}, a plataforma demonstra que os dados já
existentes e publicamente disponíveis são suficientes para identificar
desigualdades estruturais na oferta de saúde — o chamado Silêncio
Epidemiológico — e para subsidiar decisões de alocação de recursos com base
em evidências. O motor de decisão em tempo real ilustra como a tecnologia de
streaming pode transformar protocolos de triagem em sistemas inteligentes de
regulação assistencial, com potencial de redução de mortalidade evitável em
cenários de recursos limitados.

A visibilidade proporcionada pelo Silêncio Epidemiológico — municípios com
alta mortalidade e baixa oferta de leitos SUS — pode orientar futuras políticas
públicas de investimento em saúde na Região Sul, tornando o impacto do projeto
relevante além do escopo acadêmico da disciplina.


%=============================================================================
% REFERÊNCIAS BIBLIOGRÁFICAS
%=============================================================================

\begin{thebibliography}{9}

\bibitem{datasus2024}
BRASIL. Ministério da Saúde.
\textbf{DataSUS -- Transferência de Arquivos: Sistema de Informações Hospitalares
(SIH) e Cadastro Nacional de Estabelecimentos de Saúde (CNES)}.
Disponível em: \url{https://datasus.saude.gov.br/transferencia-de-arquivos/}.
Acesso em: abr. 2025.

\bibitem{ibge2023}
INSTITUTO BRASILEIRO DE GEOGRAFIA E ESTATÍSTICA (IBGE).
\textbf{Censo Demográfico 2022: Resultados do Universo}.
Rio de Janeiro: IBGE, 2023.
Disponível em: \url{https://www.ibge.gov.br/estatisticas/sociais/populacao/22827-censo-demografico-2022.html}.
Acesso em: abr. 2025.

\bibitem{lima2019}
LIMA, C. R. A. et al.
Revisão das dimensões de qualidade dos dados e métodos aplicados na avaliação
dos sistemas de informação em saúde.
\textbf{Cadernos de Saúde Pública}, Rio de Janeiro, v. 35, n. 7, 2009.
DOI: 10.1590/0102-311X00170617.

\bibitem{mendes2011}
MENDES, E. V.
\textbf{As redes de atenção à saúde}.
Brasília: Organização Pan-Americana da Saúde, 2011. 549 p.

\bibitem{spark2016}
ZAHARIA, M. et al.
Apache Spark: A unified engine for big data processing.
\textbf{Communications of the ACM}, v. 59, n. 11, p. 56--65, 2016.
DOI: 10.1145/2934664.

\bibitem{kafka2021}
APACHE SOFTWARE FOUNDATION.
\textbf{Apache Kafka: A Distributed Streaming Platform -- Documentation}.
Disponível em: \url{https://kafka.apache.org/documentation/}.
Acesso em: abr. 2025.

\bibitem{databricks2021}
DATABRICKS.
\textbf{Medallion Architecture}.
Disponível em: \url{https://www.databricks.com/glossary/medallion-architecture}.
Acesso em: abr. 2025.

\bibitem{postgresql2024}
THE POSTGRESQL GLOBAL DEVELOPMENT GROUP.
\textbf{PostgreSQL 16 Documentation}.
Disponível em: \url{https://www.postgresql.org/docs/16/}.
Acesso em: abr. 2025.

\end{thebibliography}

\end{document}
